\section{Introduction}

In classical axiomatized mathematics, "object existence" is typically established through metatheory or set theory; in semantics emphasizing computability and auditability, there is no necessary correspondence between objects that can be referenced, verified, and repeatedly re-audited and objects "declared to exist." Projective Ontological Mathematics (POM) responds to this tension via an internalized read semantics: addressability triggering read as the criterion for existence. Formally, we view the read as a partial function $\Read:\Addr\rightharpoonup \mathrm{ValDom}$; when the read is undefined, we use $\NULL$ as the notation for undefined ($\NULL\notin\mathrm{ValDom}$, and $\NULL\neq 0$).\par
Throughout this paper, we use "\textbf{fiber}" as the primary terminology (synonymous with "register/axis"; the latter two are mentioned once only at first occurrence or in fixed compounds such as prime-register).\par

\subsection{Main Results}
Under the minimal axiom system POMmin (Section \ref{sec:pommin}) and the auditable closure setting of this paper, the claims are strictly divided into two layers: \textbf{unconditional layer} (provable in the main text, cannot contain bridging assumptions) and \textbf{conditional layer} (must be presented as conditional theorem/corollary chains, providing falsifiable failure witness specifications).\par

\paragraph{Unconditional layer (POMmin + three-fiber closure; provable in main text)}
\begin{enumerate}
  \item \textbf{Partial function read and undefined semantics:} Object existence is internalized via "addressable read": exists if and only if there exists an address making the read defined and auditable; otherwise the read is undefined and denoted by $\NULL$ (Proposition \ref{prop:existence-semantics}).\par
  \item \textbf{Stable syntax and fold normalization:} We construct $\Omega_m=\{0,1\}^m$, valid subset $X_m$, and fold normalization $\Fold_m:\Omega_m\twoheadrightarrow X_m$, and provide mechanically auditable uniqueness of normal form and decidable equality fragments (Axiom \ref{ax:pommin-a3} and Appendix \ref{app:fold-rewriting}).\par
  \item \textbf{Emergence of natural number semiring:} We generate in situ operations $(\oplus,\otimes)$ on stable address space and prove their isomorphism with the semiring $(\NN,+,\times)$ (Theorem \ref{thm:emergent-N}).\par
  \item \textbf{Ledger-axis injectification and information ledger:} By externally placing fold residuals into the prime fiber (prime-register), we achieve lifting/replay and provide information ledger identities to decompose visible entropy and residual information (Proposition \ref{prop:lossless-uplift} and Theorem \ref{thm:ledger-identity}).\par
  \item \textbf{Decidability of $\mathsf{PW}_{\mathrm{dec}}$ certificate fragment:} Within the decidable fragment of projection word (PW) certificate systems, rewriting terminates and exhibits (local) confluence, yielding unique normal form; equality is decidable (Theorem \ref{thm:pw-ze-termination-confluence} and Corollary \ref{cor:pw-word-problem-decidable}).\par
  \item \textbf{Soundness under unitary slice (typed rejection):} Under the strong type auditing protocol of unitary slice, any addressable and auditable zero-point certificate must satisfy $\Re(s)=\tfrac12$; offcritical points cannot be certificated and read undefined ($\NULL$) within the closure (Theorem \ref{thm:type-safety-offcritical}).\par
  \item \textbf{Fourth fiber (register) obstacle:} Radial degrees of freedom leaving unitary slice require certificate type extension; within three-fiber closure, this requirement manifests as structural necessity of the fourth fiber (Proposition \ref{prop:fourth-register-need}).\par
\end{enumerate}

\paragraph{Conditional layer (conditional; explicit bridging premises)}
\begin{enumerate}
  \item \textbf{Address completeness $\Rightarrow$ internal $\Xi$-RH:} If address completeness assumption is adopted for fixed $L>1$ (Assumption \ref{ass:xi-address-completeness}), then all zeros of $\Xi_{\Delta,L}$ satisfy $\Re(s)=\tfrac12$ (Corollary \ref{cor:xi-rh}).\par
  \item \textbf{BridgeRH $\Rightarrow$ classical RH (conditional reduction):} If the bridging condition BridgeRH (Assumption \ref{ass:bridge-rh}) holds, then the classical RH holds (Theorem \ref{thm:bridge-rh-implies-rh}). BridgeRH is organized as a set of line-by-line auditable, falsifiable obligation lists: each obligation should provide either a finite certificate or a failure witness (see Appendix \ref{app:bridge-rh-obligations}).\par
\end{enumerate}

\subsection{Related Work and Paper Positioning}
This paper's three main threads correspond to three directions in existing literature, unified here as "auditable read—ledger-axis injectification—precision ledger—typed rejection" in the typed mathematics layer.\par
\textbf{(1) Stable syntax and fold normalization:} The golden ratio syntax and its inverse limit structure can be viewed as SFT (shift of finite type) in symbolic dynamics; its counting and coding properties can be understood within the framework of symbolic dynamics and coding theory \cite{LindMarcus1995}. Fold normalization employs terminating/(locally) confluent rewriting systems to ensure unique normal form, consistent with combinatorial rewriting theory and can be organized via standard tools such as Newman's lemma for auditable equality chains \cite{Newman1942}.\par
\textbf{(2) Ledger-axis injectification and information ledger:} The information ledger identity employs standard decomposition of Shannon entropy and KL divergence \cite{CoverThomas2006}, decomposing the irreversibility caused by many-to-one fold into "fiber volume terms" and "residual terms offset from maximum entropy uplift"; the residual bearing of prime-register and Gödel-like coding organize the discrete trajectory information needed for replay as finite verifiable data.\par
\textbf{(3) Auditing of unitary slice, proof-carrying, and continuous precision ledger:} The unitary slice translates critical line geometry into the unit circle parameter domain (Proposition \ref{prop:critical-line-unitary-slice}) and can be further algebraized to reciprocal polynomials and interval root criteria; the related unit circle spectral structure and numerical algorithms have mature frameworks in OPUC/CMV and spectral theory \cite{CanteroMoralVelazquez2003CMV,Simon2005OPUC1,BevilacquaDelCorsoGemignani2015CMVCompanion}. We emphasize its "proof-carrying" version: employing outward-rounding interval arithmetic and rational certificates (e.g., Schur--Cohn/LMI and WSOS) as offline re-computable auditing backend \cite{DaumasLesterMunoz2007,DymYoung1990SchurCohnMatrixPolynomials,DavisPapp2023}. On the other hand, unit circle area phase compactification compresses unbounded scales to boundary phase, and via Jacobian-induced precision potential, transcodes "tail compression cost" to auditable precision ledger (see Appendices \ref{app:phase-compactification-register-unification} and \ref{app:structural-synthesis-phase-diagram}).\par
The role of BridgeRH is to organize the zero events of classical completion functions into addressified reads within the auditable object domain; its form is compatible with the trace--primitive--Euler strict exchange bridge structure of $\zeta$ functions in dynamical systems \cite{Ruelle1976ZetaExpanding,Ruelle1978Thermodynamic,ParryPollicott1990Zeta}. We do not attempt to provide classical analytic number theoretic proofs at the same level as the main text (see standard references \cite{Apostol1976}), but rather explicitly decompose the difficulty into line-by-line auditable bridge obligations, marking its minimal failure points with $\NULL$ semantics and fourth-register obstacle.\par

\subsection{Novelty audit: contribution inventory and boundary declarations}
For the convenience of review and verification, we classify the material into three categories and annotate the logical status of each.\par
\begin{itemize}
  \item \textbf{Standard tools (demoted to background/lemma):} Zeckendorf/Ostrowski representations and inverse limits, Newman's lemma and combinatorial rewriting, Shannon entropy and $D_{\mathrm{KL}}$ decomposition, interval arithmetic and outward rounding, etc.; we employ these only as organizational tools for mechanically auditable chains.\par
  \item \textbf{Reorganized but independently publishable structures:} Partial function read (certificate existence semantics) and $\NULL$ undefined discipline; conservative expression of fold—residual—replay and auditing interface; typed rejection mechanism of unitary slice and soundness theorem.\par
  \item \textbf{Contributions that must be highlighted and strengthened:} Completeness of record structure (canonical four-factor generators of minimal ledger-axis; details in Appendix \ref{app:structural-synthesis-phase-diagram}); fourth-fiber obstacle's structure and quantitative dual characterization (type extension and precision load lower bounds; see Section \ref{sec:fourth-register} and related appendices); falsifiable specification of BridgeRH system (each assumption must be equipped with failure witness and can terminate corollary chains; see Section \ref{sec:bridge-rh} and Appendix \ref{app:bridge-rh-obligations}).\par
\end{itemize}

\subsection{Paper structure}
Section \ref{sec:pommin} provides minimal axioms of POMmin and existence semantics. We then provide scan—projection generators and three-fiber structure (also called three-registers: scale/phase/prime fiber), and establish arithmetic emergence (Section \ref{sec:emergent-arithmetic}) and information ledger identities (Section \ref{sec:ledger}) on stable address space. We then discuss projection word certificate semantics, unitary slice constraints and type safety, and provide conditional inference from internal $\Xi$-statements to classical RH under bridging assumptions (Section \ref{sec:bridge-rh}); the fourth-fiber (register) obstacle is in Section \ref{sec:fourth-register}. Section \ref{sec:conclusion-cluster} provides certificates and unified coding layer (self-describing codes, monocode residuals, and repeated auditability of unitary slice), and provides support for verifiable conditions of BridgeRH; discussion and conclusion appear in subsequent sections, with appendices collecting proof-carrying constructions and related technical materials.\par
